
\documentclass{article}
\usepackage{colortbl}
\usepackage{makecell}
\usepackage{multirow}
\usepackage{supertabular}

\begin{document}

\newcounter{utterance}

\centering \large Interaction Transcript for game `ifeval', experiment `ifeval', episode 161 with reflexion{-}llama{-}3.1{-}8b{-}instruct.
\vspace{24pt}

{ \footnotesize  \setcounter{utterance}{1}
\setlength{\tabcolsep}{0pt}
\begin{supertabular}{c@{$\;$}|p{.15\linewidth}@{}p{.15\linewidth}p{.15\linewidth}p{.15\linewidth}p{.15\linewidth}p{.15\linewidth}}
    \# & \multicolumn{2}{c}{Player} && \multicolumn{2}{c}{Game Master} \\
    \hline

    \theutterance \stepcounter{utterance}  
    & & & \multicolumn{4}{p{0.6\linewidth}}{
        \cellcolor[rgb]{0.9,0.9,0.9}{
            \makecell[{{p{\linewidth}}}]{
                \texttt{\tiny{[P1$\langle$GM]}}
                \texttt{You have attempted this same task before and did not succeed. These are the plans you wrote after those attempts.} \\
\texttt{Trial \#1: New plan: } \\
\\ 
\texttt{Upon reflection, I realize that my previous attempt was too focused on describing the performance itself, without considering the specific task of continuing and expanding the given sentence. I got caught up in the excitement of the performance and lost sight of the goal.} \\
\\ 
\texttt{To improve, I will take a step back and re{-}examine the original sentence: "A man is performing for a crowd on a stage, ..." My task is to continue and expand this sentence, not to create a new story or describe a different scene.} \\
\\ 
\texttt{Next time, I will focus on building on the existing sentence, using the given information as a starting point. I will ask myself questions like: What is the man doing on stage? What is the crowd's reaction? What is the atmosphere like?} \\
\\ 
\texttt{To achieve this, I will take a more structured approach, breaking down the task into smaller, manageable chunks. I will start by identifying the key elements of the original sentence, such as the man, the crowd, and the stage. Then, I will use these elements as a foundation to build on, adding new details and descriptions to continue and expand the sentence.} \\
\\ 
\texttt{For example, I might start by describing the man's actions on stage, such as singing or playing an instrument. I might then describe the crowd's reaction, including their emotions and physical responses. Finally, I might add more details about the stage and the atmosphere, using sensory language to bring the scene to life.} \\
\\ 
\texttt{By taking a more focused and structured approach, I hope to create a more cohesive and engaging continuation of the original sentence.} \\
\\ 
\texttt{Trial \#2: New plan: } \\
\\ 
\texttt{Upon reflection, I realize that my previous attempt was too focused on describing the performance itself, without considering the specific task of continuing and expanding the given sentence. I got caught up in the excitement of the performance and lost sight of the goal.} \\
\\ 
\texttt{To improve, I will take a step back and re{-}examine the original sentence: "A man is performing for a crowd on a stage, ..." My task is to continue and expand this sentence, not to create a new story or describe a different scene.} \\
\\ 
\texttt{Next time, I will focus on building on the existing sentence, using the given information as a starting point. I will ask myself questions like: What is the man doing on stage? What is the crowd's reaction? What is the atmosphere like?} \\
\\ 
\texttt{To achieve this, I will take a more structured approach, breaking down the task into smaller, manageable chunks. I will start by identifying the key elements of the original sentence, such as the man, the crowd, and the stage. Then, I will use these elements as a foundation to build on, adding new details and descriptions to continue and expand the sentence.} \\
\\ 
\texttt{However, I also realize that my previous attempt was too repetitive and lacked variety. I repeated the same descriptions and ideas multiple times, without adding any new insights or perspectives.} \\
\\ 
\texttt{To address this, I will make a conscious effort to vary my descriptions and add more depth and nuance to the scene. I will use sensory language to bring the scene to life, incorporating details about the sights, sounds, and emotions of the performance.} \\
\\ 
\texttt{For example, I might describe the man's voice as "rich and velvety," or the crowd's reaction as "electric and frenzied." I might also add more details about the stage and the atmosphere, such as the smell of sweat and smoke, or the feeling of the crowd's energy.} \\
\\ 
\texttt{By taking a more focused and structured approach, and varying my descriptions to add more depth and nuance, I hope to create a more engaging and immersive continuation of the original sentence.} \\
\\ 
\\ 
\texttt{Can you please continue and expand: "A man is performing for a crowd on a stage, ..."} \\
\texttt{Use at least 60 sentences, but less than 600 words.} \\
            }
        }
    }
     \\ \\

    \theutterance \stepcounter{utterance}  
    & \multicolumn{4}{p{0.6\linewidth}}{
        \cellcolor[rgb]{0.561,0.737,0.561}{
            \makecell[{{p{\linewidth}}}]{
                \texttt{\tiny{[P1$\rangle$GM]}}
                \texttt{A man is performing for a crowd on a stage, his voice soaring through the microphone as he belts out a soulful melody. The crowd is entranced, their faces aglow with a mix of wonder and adoration. They sway to the rhythm, their bodies moving in perfect sync with the music.} \\
\\ 
\texttt{The man's voice is rich and velvety, like a warm summer night that wraps itself around the listener's heart. It's a voice that can convey the depths of human emotion, from the heights of joy to the depths of sorrow. As he sings, the crowd is transported to a different world, one where the cares and worries of everyday life are left behind.} \\
\\ 
\texttt{The stage is a sea of lights and colors, a kaleidoscope of hues that dance and swirl around the performer. The crowd is bathed in a warm, golden glow, as if the very essence of the music has been distilled into a physical presence. The air is electric, charged with the energy of the performance.} \\
\\ 
\texttt{As the man sings, his voice grows stronger, more confident. He's a master of his craft, a virtuoso who weaves a spell of enchantment over the crowd. They're under his spell, helpless to resist the power of his music. They're swept up in a tide of emotion, their hearts pounding in time with the beat.} \\
\\ 
\texttt{The crowd is a sea of faces, a diverse tapestry of humanity that's been woven together by the shared experience of the performance. There are young and old, rich and poor, black and white, all united in their love of music. They're a community, a tribe that's been brought together by the power of art.} \\
\\ 
\texttt{As the man sings, the crowd begins to respond. They clap and cheer, their voices rising in a crescendo of approval. They're caught up in the moment, lost in the beauty and wonder of the performance. They're alive, connected to something deeper and more meaningful than themselves.} \\
\\ 
\texttt{The stage is a platform, a stage upon which the performer stands tall and proud. It's a symbol of the power of art, a reminder that music has the ability to transcend borders and boundaries. It's a language that speaks to the heart, a universal tongue that's understood by all.} \\
\\ 
\texttt{As the man sings, the crowd is transported to a different time and place. They're in a world of their own, a world that's been created by the power of music. They're in a world of dreams and imagination, a world where anything is possible.} \\
\\ 
\texttt{The man's voice is a instrument, a tool that's been crafted to convey the depths of human emotion. It's a voice that can make you laugh and cry, a voice that can make you feel alive. As he sings, the crowd is reminded of the power of music, the ability it has to touch the heart and soul.} \\
\\ 
\texttt{The crowd is a living, breathing entity, a being that's been brought to life by the power of the performance. They're a community, a tribe that's been united by the shared experience of the music. They're a family, a group of people who are connected by their love of art.} \\
\\ 
\texttt{As the man sings, the crowd begins to respond. They clap and cheer, their voices rising in a crescendo of approval. They're caught up in the moment, lost in the beauty and wonder of the performance. They're alive, connected to something deeper and more meaningful than themselves.} \\
\\ 
\texttt{The stage is a platform, a stage upon which the performer stands tall and proud. It's a symbol of the power of art, a reminder that music has the ability to transcend borders and boundaries. It's a language that speaks to the heart, a universal tongue that's understood by all.} \\
\\ 
\texttt{As the man sings, the crowd is transported to a different time and place. They're in a world of their own, a world that's been created by the power of music. They're in a world of dreams and imagination, a world where anything is possible.} \\
\\ 
\texttt{The man's voice is a instrument, a tool that's been crafted to convey the depths of human emotion. It's a voice that can make you laugh and cry, a voice that can make you feel alive. As he sings, the crowd is reminded of the power of music, the ability it has to touch the heart and soul.} \\
\\ 
\texttt{The crowd is a living, breathing entity, a being that's been brought to life by the power of the performance. They're a community, a tribe that's been united by the shared experience of the music. They're a family, a group of people who are connected by their love of art.} \\
\\ 
\texttt{As the man sings, the crowd begins to respond. They clap and cheer, their voices rising in a crescendo of approval. They're caught up in the moment, lost in the beauty and wonder of the performance. They're alive, connected to something deeper and more meaningful than themselves.} \\
\\ 
\texttt{The stage is a platform, a stage upon which the performer stands tall and proud. It's a symbol of the power of art, a reminder that music has the ability to transcend borders and boundaries. It's a language that speaks to the heart, a universal tongue that's understood by all.} \\
\\ 
\texttt{As the man sings, the crowd is transported to a different time and place. They're in a world of their own, a world that's been created by the power of music. They're in a world of dreams and imagination, a world where anything is possible.} \\
\\ 
\texttt{The man's voice is a instrument, a tool that's been crafted to convey the depths of human emotion. It's a voice that can make you laugh and cry, a voice that can make you feel alive. As he sings, the crowd is reminded of the power of music, the ability it has to touch the heart and soul.} \\
\\ 
\texttt{The crowd is a living, breathing entity, a being that's been brought to life by the power of the performance. They're a community, a tribe that's been united by the shared experience of the music. They're a family, a group of people who are connected by their love of art.} \\
\\ 
\texttt{As the man sings, the crowd begins to respond. They clap and cheer, their voices rising in a crescendo of approval. They're caught up in the moment, lost in the beauty and wonder of} \\
            }
        }
    }
    & & \\ \\

    \theutterance \stepcounter{utterance}  
    & & & \multicolumn{2}{p{0.3\linewidth}}{
        \cellcolor[rgb]{0.9,0.9,0.9}{
            \makecell[{{p{\linewidth}}}]{
                \texttt{\tiny{[GM$|$GM]}}
                {'length_constraints:number_sentences': {'relation': 'at least', 'num_sentences': 60}, 'length_constraints:number_words': {'relation': 'less than', 'num_words': 600}}
            }
        }
    }
    & & \\ \\

    \theutterance \stepcounter{utterance}  
    & & & \multicolumn{2}{p{0.3\linewidth}}{
        \cellcolor[rgb]{0.9,0.9,0.9}{
            \makecell[{{p{\linewidth}}}]{
                \texttt{\tiny{[GM$|$GM]}}
                \texttt{game\_result = LOSE} \\
            }
        }
    }
    & & \\ \\

    \theutterance \stepcounter{utterance}  
    & & \multicolumn{4}{p{0.6\linewidth}}{
        \cellcolor[rgb]{0.95,0.95,0.95}{
            \makecell[{{p{\linewidth}}}]{
                \texttt{\tiny{[P1$\leftrightarrow$P]}}
                \texttt{Post{-}episode Reflexion (not sent to GM):} \\
\texttt{New plan: } \\
\\ 
\texttt{Upon reflection, I realize that my previous attempts were too focused on describing the performance itself, without considering the specific task of continuing and expanding the given sentence. I got caught up in the excitement of the performance and lost sight of the goal.} \\
\\ 
\texttt{To improve, I will take a step back and re{-}examine the original sentence: "A man is performing for a crowd on a stage, ..." My task is to continue and expand this sentence, not to create a new story or describe a different scene.} \\
\\ 
\texttt{Next time, I will focus on building on the existing sentence, using the given information as a starting point. I will ask myself questions like: What is the man doing on stage? What is the crowd's reaction? What is the atmosphere like?} \\
\\ 
\texttt{To achieve this, I will take a more structured approach, breaking down the task into smaller, manageable chunks. I will start by identifying the key elements of the original sentence, such as the man, the crowd, and the stage. Then, I will use these elements as a foundation to build on, adding new details and descriptions to continue and expand the sentence.} \\
\\ 
\texttt{However, I also realize that my previous attempts were too repetitive and lacked variety. I repeated the same descriptions and ideas multiple times, without adding any new insights or perspectives.} \\
\\ 
\texttt{To address this, I will make a conscious effort to vary my descriptions and add more depth and nuance to the scene. I will use sensory language to bring the scene to life, incorporating details about the sights, sounds, and emotions of the performance.} \\
\\ 
\texttt{Additionally, I will pay closer attention to the sentence structure and grammar, making sure to use proper punctuation and sentence flow to create a cohesive and engaging narrative.} \\
\\ 
\texttt{To achieve this, I will take a more focused and structured approach, varying my descriptions to add more depth and nuance, and paying closer attention to sentence structure and grammar. I will also make a conscious effort to add more variety and creativity to my descriptions, using sensory language and vivid imagery to bring the scene to life.} \\
\\ 
\texttt{For example, I might start by describing the man's actions on stage, such as singing or playing an instrument. I might then describe the crowd's reaction, including their emotions and physical responses. Finally, I might add more details about the stage and the atmosphere, using sensory language to bring the scene to life.} \\
\\ 
\texttt{By taking a more focused and structured approach, varying my descriptions to add more depth and nuance, and paying closer attention to sentence structure and grammar, I hope to create a more engaging and immersive continuation of the original sentence.} \\
            }
        }
    }
    & \\ \\

\end{supertabular}
}

\end{document}
